\chapter{Integrál}
\begin{chapquote}{Isaac Asimov}
Po letech, kdy mi matematika připadala snadná, jsem se konečně dostal k integrálnímu počtu a narazil jsem na překážku. Uvědomil jsem si, že dál už se nedostanu, a dodnes se mi nepodařilo tuto hranici překročit.
\end{chapquote}
Doteď jsme řešili diferenciální rovnice tak, že jsme řešení buď uhodli, nebo jsme si vymysleli novou funkci a tu jsme nazvali řešením. V příštích dvou kapitolách se naučíme, jak (některé) rovnice řešit systematicky. Zkonstruujeme operaci, která umí v diferenciální rovnici \uv{odstranit} derivaci. Této operaci se říká \emph{integrál}.
\section{Plocha pod grafem}
\section{Primitivní funkce}
\section{Metoda separace proměnných}%logisticka krivka, mv'=-g-kv^2, vytekani vody, chemicka kinetika
